%%
%% spring-lecture01.tex
%% 
%% Made by Alex Nelson <pqnelson@gmail.com>
%% Login   <alex@lisp>
%% 
%% Started on  2026-04-09T09:29:10-0700
%% Last update 2026-04-09T09:29:10-0700
%% 

\lecture{}

\begin{node}
This term will be about Riemannian geometry.
\end{node}

\begin{node}
What is Riemannian geometry? The study of spaces where we can compute
lengths and angles. We will blast through a review (of smooth
manifolds) quickly. We equip a manifold with a Riemannian
structure. Using length and angles, we can introduce a notion of curvature.
\end{node}

\subsection{Topological and smooth manifolds}

\begin{definition}[Topological manifold]
A \define{Topological Manifold} consists of a Hausdorff,
second-countable topological space $M$ equipped with an open cover
$\{U_{\alpha}\}_{\alpha\in A}$ and a set of continuous maps
$\varphi_{\alpha}\colon U_{\alpha}\to\RR^{m}$ which are homeomorphisms
into their image $\varphi_{\alpha}(U_{\alpha})$ which is an open
subset of $\RR^{m}$. (The $m$ is fixed for $M$ and it will not vary
chart to chart). Here the $m$ is the \define{Dimension} of $M$ (as a
topological manifold) and we will say that $M$ is a ``topological $m$-manifold''.
\end{definition}

\begin{definition}
We call the collection $\{(U_{\alpha},\varphi_{\alpha})\}_{\alpha\in A}$
a \define{Topological Atlas} on $M$.

We call an individual pair $(U_{\alpha},\varphi_{\alpha})$ a
\define{Local Chart}.
\end{definition}

\begin{node}
We define smooth manifolds by borrowing smoothness from $\RR^{n}$.
\end{node}

\begin{definition}[Smooth atlas]
Let $M$ be a topological manifold, let
$\mathcal{A}=\{(U_{\alpha},\varphi_{\alpha}\}_{\alpha\in A}$ be a
topological atlas on $M$. We say the atlas $\mathcal{A}$ is
\define{Smooth} if for every $\alpha,\beta\in A$ such
that\footnote{This non-disjointness of patches condition is, strictly
speaking, un-necessary because the only possible transition map would
be the identity map, which is always arbitrarily smooth.}
$U_{\alpha}\cap U_{\beta}\neq\emptyset$ we have
\begin{equation}
\varphi_{\beta}\circ\varphi_{\alpha}^{-1}\colon\varphi_{\alpha}(U_{\alpha}\cap U_{\beta})\to\varphi_{\beta}(U_{\alpha}\cap U_{\beta})
\end{equation}
is smooth in the sense of calculus (since it is a function between two
open subsets of $\RR^{n}$).

We call the functions $\varphi_{\beta}\circ\varphi_{\alpha}^{-1}$ the
\define{Transition Functions} of the atlas $\mathcal{A}$.
\end{definition}

\begin{definition}[Smooth manifold]
A \define{Smooth $n$-Manifold} is a topological $n$-manifold $M$
equipped with one smooth atlas.

We can speak of other regularity classes: replace ``smooth'' [i.e.,
  $C^{\infty}$] with ``$C^{k}$'' everywhere to get $C^{k}$ atlases and
$C^{k}$ manifolds. When $k=0$, we see $C^{0}$ manifolds are just
topological manifolds.
\end{definition}

\begin{node}
Is every topological manifold a smooth manifold? No! This was proven
in the 1980s by Donaldson and Simon. There are 4-dimensional examples
of such spaces.
\end{node}

\begin{definition}[Compatibility]
Let $M$ be a smooth manifold with a smooth atlas with a smooth atlas
$\mathcal{A}=\{(U_{\alpha},\varphi_{\alpha})\}_{\alpha\in A}$. Let
$\varphi\colon U\to\RR^{n}$ be another chart of $M$. We say
$(U,\varphi)$ is \define{Compatible} with the atlas $\mathcal{A}$ if
for each $\alpha\in A$ such that $U_{\alpha}\cap U\neq\emptyset$ we
have the transition function $\varphi\circ\varphi_{\alpha}^{-1}\colon\varphi_{\alpha}(U_{\alpha}\cap U)\to\varphi(U_{\alpha}\cap U)$
is smooth.

Again, we can replace ``smooth'' with ``$C^{k}$''.
\end{definition}

\begin{node}
If a chart $(U,\varphi)$ is compatible with an atlas $\mathcal{A}$,
then we can form an enlarged atlas $\mathcal{A}'=\mathcal{A}\cup\{(U,\varphi)\}$.
\end{node}

\begin{remark}
This is actually a little too slick. We can speak of two charts being
$C^{k}$ compatible (for $k\in\NN_{0}\cup\{\infty,\omega\}$), which is
an equivalence relation. Then we can speak of two $C^{k}$ atlases
$\mathcal{A}_{1}$ and $\mathcal{A}_{2}$ being $C^{k}$-equivalent if all their
charts are compatible (or equivalently, if
$\mathcal{A}_{1}\cup\mathcal{A}_{2}$ forms a $C^{k}$ atlas). This
gives us an equivalence relation of $C^{k}$ atlases. Moreover, the
collection of all possible $C^{k}$ atlases forms a poset under
inclusion, and we can use Zorn's lemma to get the maximal atlas for
any given $C^{k}$ atlas.
\end{remark}

\begin{definition}[Smooth structure]
Let $M$ be a smooth manifold. The \define{Smooth Structure} of $M$ is
the unique maximal smooth atlas containing the one you endowed $M$
with.

(Technically, this is the $C^{\infty}$ differential structure of
$M$. You could consider different $C^{k}$ differential structures, for $k>0$.)
\end{definition}

\begin{notation}
We write $\smoothstr{M}$ for the smooth structure of $M$.
\end{notation}

\begin{remark}
Every topological 2-manifold, and every topological 3-manifold, admits
a unique smooth structure.
\end{remark}

\begin{remark}
There are topological 4-manifolds that are \emph{non-smoothable}
(i.e., cannot be endowed with any smooth atlas). This is due to
Donaldson's work from the 1980s.
\end{remark}

\begin{remark}
There are uncountably many distinct smooth structures on $\RR^{4}$, a
result due to Taubes from the 1980s.
\end{remark}

\begin{remark}
\textsc{Smooth Poincar\'{e} Conjecture} remains an open problem: How
many distinct smooth structures are there on $\sphere{4}$?
\end{remark}

\begin{definition}[Embedded smooth manifolds of $\RR^{k}$]
A subset $M\subset\RR^{k}$ is called a \define{Smooth $n$-Submanifold of $\RR^{k}$}
if for every $x\in M$ there exists an open subset $V\subset\RR^{k}$,
an open subset $U\subset M$, and a smooth $\varphi\colon U\to\RR^{k}$
such that
\begin{enumerate}
\item $\varphi\colon U\to M\cap V$ is a homeomorphism
\item $\D\varphi(y)\colon\RR^{n}\to\RR^{k}$ is injective for each
  $y\in U$.
\end{enumerate}
When $k=n+1$, then $M$ is called a \define{Hypersurface} of $\RR^{k}$.
\end{definition}

\begin{proposition}
Every smooth $n$-submanifold of $\RR^{k}$ is a smooth $n$-manifold.
\end{proposition}

This is a consequence of either the implicit function theorem, or the
inverse function theorem.

\begin{node}
So the induced topology of $M$ is automatically Hausdorff and
second-countable. The charts are formed from $\varphi\colon U\to M\cap V$
by taking $(M\cap V,\varphi^{-1})$. These are all smooth.
\end{node}

\begin{example}
Let $f\colon U\subset\RR^{n}\to\RR$ be a smooth function. Then the
graph of $f$,
\begin{equation}
\Gamma(f)=\{(x_{1},\dots,x_{n},f(x_{1},\dots,x_{n}))\in U\times\RR\}\subset\RR^{n+1}
\end{equation}
is a smooth submanifold of $\RR^{n+1}$ always coverable by a single chart.
\end{example}

\endinput